%! Unit 5: Radical Functions — Summative Assessment — Unit Test (blank)
\documentclass[10pt]{article}
\usepackage{algebra2-article}
\usepackage{algebra2-boxes}
% Coordinate grid: {xmin}{xmax}{ymin}{ymax}
\newcommand{\lgrid}[4]{%
  \draw[step=1, linegray!40, very thin] (#1,#3) grid (#2,#4);
  \draw[->, semithick, charcoal] (#1,0)--(#2,0) node[below right, font=\scriptsize]{$x$};
  \draw[->, semithick, charcoal] (0,#3)--(0,#4) node[left, font=\scriptsize]{$y$};}
\newcommand{\dpt}[2]{\fill[charcoal] (#1,#2) circle (2.6pt);}

% Part-divider strip (headlinebox takes one arg: the background color)
\newcommand{\parthead}[1]{%
  \vspace{6pt}\noindent
  \begin{headlinebox}{forest}\color{white}\bfseries #1\end{headlinebox}%
  \vspace{4pt}}

\begin{document}

\pageheader{Unit 5: Radical Functions}{Unit Test}
\namedateperiod

\vspace{0.06in}
\noindent\textbf{Instructions:} Show all work. No notes unless your teacher says so.

% ======================================================================
\parthead{Part A --- Vocabulary (8 pts)}
Write the letter of the best description next to each term.

\vspace{4pt}
\noindent
\begin{minipage}[t]{0.36\textwidth}
\begin{enumerate}[leftmargin=2.2em, itemsep=7pt, label=\arabic*.]
  \item Index \blank{1.2cm}
  \item Radicand \blank{1.2cm}
  \item Principal root \blank{1.2cm}
  \item Conjugate \blank{1.2cm}
  \item Extraneous solution \blank{1.2cm}
  \item Composition \blank{1.2cm}
  \item Inverse function \blank{1.2cm}
  \item Anchor point \blank{1.2cm}
\end{enumerate}
\end{minipage}%
\hfill
\begin{minipage}[t]{0.61\textwidth}
\small
\begin{description}[leftmargin=1.4em, itemsep=3pt]
  \item[(A)] A candidate that satisfies the powered-up equation but fails when it is tested in the original.
  \item[(B)] The non-negative value the radical symbol by itself returns; a $\pm$ must be written in by hand.
  \item[(C)] The function that reverses another's inputs and outputs; its graph is the mirror image over the line $y=x$.
  \item[(D)] A binomial with the same two terms but the opposite middle sign; multiplying by it turns an irrational denominator rational.
  \item[(E)] The small number in the crook of the radical symbol that says which power is being undone.
  \item[(F)] The point $(h,k)$ that positions a transformed radical graph --- an endpoint when the index is even, a turning point when it is odd.
  \item[(G)] Everything written underneath the radical symbol; when the index is even, it may not be negative.
  \item[(H)] The operation that feeds one function's output into another function's input, written $(f\circ g)(x)$.
\end{description}
\end{minipage}

% ======================================================================
\parthead{Part B --- Multiple Choice (12 pts)}
Circle the best answer.

\begin{enumerate}[leftmargin=1.9em, itemsep=9pt]
  \item Written with a rational exponent, $\sqrt[5]{x^4}$ is
    \hfill (A) $x^{5/4}$ \quad (B) $x^{4/5}$ \quad (C) $x^{20}$ \quad (D) $x^{-5/4}$
  \item The domain of $f(x)=\sqrt{9-3x}$ is
    \hfill (A) $x\le 3$ \quad (B) $x\ge 3$ \quad (C) $x\ge -3$ \quad (D) all real numbers
  \item The graph of $y=\sqrt{2x-10}+3$ begins at the point
    \hfill (A) $(10,3)$ \quad (B) $(5,3)$ \quad (C) $(5,-3)$ \quad (D) $(-5,3)$
  \item In simplest form, $\sqrt{98}+\sqrt{32}$ is
    \hfill (A) $\sqrt{130}$ \quad (B) $28\sqrt2$ \quad (C) $11\sqrt{130}$ \quad (D) $11\sqrt2$
  \item Which equation has \emph{no real solution}?
    \\[3pt]\hspace*{1.2em} (A) $\sqrt[3]{x+2}=-3$
    \hspace{1.0cm} (B) $\sqrt{x+1}=-5$
    \hspace{1.0cm} (C) $\sqrt{x-4}=6$
    \hspace{1.0cm} (D) $\sqrt[3]{x}=0$
  \item If $f(x)=\sqrt{x}$ and $g(x)=x+4$, then $(f\circ g)(x)=$
    \hfill (A) $\sqrt{x+4}$ \quad (B) $\sqrt{x}+4$ \quad (C) $\sqrt{4x}$ \quad (D) $x\sqrt{x}+4\sqrt{x}$
\end{enumerate}

% ======================================================================
\parthead{Part C --- Short Answer \& Computation (40 pts)}
Show your work. Assume every variable is non-negative unless told otherwise.

\begin{enumerate}[leftmargin=1.9em, itemsep=10pt]
  \item \textbf{Read the graph.} The radical function $f$ is graphed below; each gridline is one unit
        and the solid dot is the endpoint.
        \begin{center}
        \begin{tikzpicture}[scale=0.45]
          \lgrid{-6}{8}{-4}{4}
          \begin{scope}
            \clip (-6,-4) rectangle (8,4);
            \draw[very thick, forest, domain=-4:8, samples=120, smooth]
              plot (\x,{sqrt((\x)+4)-1});
          \end{scope}
          \fill[forest] (-4,-1) circle (4pt);
          \foreach \x in {-5,-2,2,4,6}
            \draw[charcoal] (\x,-0.15)--(\x,0.15) node[below=1pt, font=\tiny]{$\x$};
          \dpt{-3}{0}
          \dpt{0}{1}
          \dpt{5}{2}
        \end{tikzpicture}
        \end{center}
        (a) Coordinates of the endpoint: \blank{2.4cm}
        \\[4pt]
        (b) Domain (interval notation): \blank{3.0cm}; \quad range: \blank{3.0cm}
        \\[4pt]
        (c) $x$-intercept: \blank{2.0cm}; \quad $y$-intercept: \blank{2.0cm}
        \\[4pt]
        (d) Interval(s) where $f$ is increasing: \blank{4.0cm}
        \\[4pt]
        (e) Does $f$ have an absolute maximum or an absolute minimum? Give it and say where it occurs.
        \\[2pt] \blank{11.0cm}
        \\[4pt]
        (f) End behavior: as $x\to\infty$, $f(x)\to$ \blank{2.0cm};
            \ as $x\to-\infty$, \blank{4.6cm}
        \\[4pt]
        (g) Interval(s) where $f(x)<0$: \blank{4.0cm}

  \item \textbf{Radicals and rational exponents.} Rewrite or evaluate. Leave no negative exponents.
        \\[4pt]
        (a) Write $\sqrt[3]{x^2}$ with a rational exponent. \blank{3.0cm}
        \\[4pt]
        (b) Write $y^{7/2}$ as a radical. \blank{3.0cm}
        \\[4pt]
        (c) $8^{2/3}=$ \blank{2.2cm} \hspace{1.2cm}
        (d) $25^{-1/2}=$ \blank{2.2cm}
        \\[4pt]
        (e) Write $x^{3/4}\cdot x^{1/2}$ as a single radical.
        \vspace{1.2cm}

  \item \textbf{Simplify. Add or subtract where you can.}
        \\[6pt]
        (a) $\sqrt{98}$ \hspace{2.4cm} (b) $\sqrt[3]{40}$
        \vspace{1.9cm}
        \\
        (c) $\sqrt{18x^5}$ \hspace{2.1cm} (d) $3\sqrt{8}-\sqrt{50}$
        \vspace{2.1cm}

  \item \textbf{Multiply, then rationalize.} Every answer must be in simplest form.
        \\[6pt]
        (a) $\sqrt{10}\cdot\sqrt{14}$ \hspace{1.8cm} (b) $\left(3+\sqrt7\right)\left(3-\sqrt7\right)$
        \vspace{2.1cm}
        \\
        (c) $\dfrac{10}{\sqrt5}$ \hspace{3.0cm} (d) $\dfrac{6}{4-\sqrt{10}}$
        \vspace{2.5cm}

  \item \textbf{Transformations.} Let $g(x)=-3\sqrt{x-2}+1$.
        \\[4pt]
        (a) List every transformation that takes $y=\sqrt{x}$ to $g$. \blank{6.6cm}
        \\[4pt]
        (b) Anchor point: \blank{2.2cm}
        \\[4pt]
        (c) Domain: \blank{2.8cm}; \quad range: \blank{2.8cm}
        \\[4pt]
        (d) Is the anchor an absolute maximum or an absolute minimum? \blank{4.0cm}
        \\[4pt]
        (e) A different graph of the form $y=\sqrt{x-h}+k$ begins at $(-4,3)$. Write its equation.
            \blank{3.6cm}
        \\[4pt]
        (f) For $p(x)=\sqrt[3]{x+5}-1$, give the anchor point \blank{2.2cm} and the domain
            \blank{2.6cm}. In one sentence, say why that domain is not like $g$'s.
        \\[2pt] \blank{12.0cm}

  \item \textbf{Solve. Test every candidate in the original equation.}
        \\[6pt]
        (a) $\sqrt{3x+4}=x$
        \vspace{2.9cm}
        \\
        (b) $\sqrt[3]{x-1}=-2$
        \vspace{2.1cm}
        \\
        (c) $\sqrt{2x-5}+7=3$
        \vspace{2.1cm}
        \\
        (d) $x^{2/3}=16$
        \vspace{2.1cm}

  \item \textbf{Composition.} Let $f(x)=\sqrt{x}$ and $g(x)=x-9$.
        \\[4pt]
        (a) $(f\circ g)(x)=$ \blank{3.0cm}, \quad domain: \blank{2.6cm}
        \\[4pt]
        (b) $(g\circ f)(x)=$ \blank{3.0cm}, \quad domain: \blank{2.6cm}
        \\[4pt]
        (c) Are these the same function? Answer using \emph{both} the rule and the domain.
        \\[2pt] \blank{12.0cm}
        \\[4pt]
        (d) $(f\circ g)(25)=$ \blank{2.2cm}
        \\[4pt]
        (e) For $p(x)=3x-2$ and $q(x)=x^2+1$, find $(p\circ q)(-3)$. \blank{3.0cm}

  \item \textbf{Inverses.}
        \\[4pt]
        (a) $f(x)=5x+3$. Find $f^{-1}(x)$. \blank{4.0cm}
        \\[4pt]
        (b) Verify by composition --- \emph{both} orders --- that your answer really is the inverse.
        \vspace{2.6cm}
        \\
        (c) $g(x)=\sqrt{x-4}$. Find $g^{-1}(x)$ and state its domain.
        \vspace{2.4cm}
        \\
        (d) The point $(13,3)$ lies on the graph of $g$. What point must lie on the graph of
            $g^{-1}$? \blank{2.6cm}
\end{enumerate}

% ======================================================================
\parthead{Part D --- Extended Response (12 pts)}
Explain your reasoning in full sentences.

\begin{enumerate}[leftmargin=1.9em, itemsep=16pt]
  \item \textbf{The full analysis.} Consider $f(x)=\sqrt{x-4}+1$.
        \textbf{(a)} Give the anchor point, the domain, and the range, both in interval notation.
        \textbf{(b)} Find the $x$-intercept and the $y$-intercept, or explain why one does not exist.
        \textbf{(c)} List the transformations that take $y=\sqrt{x}$ to $f$.
        \textbf{(d)} Find $f^{-1}(x)$ and state its domain, saying exactly where that restriction
        comes from.
        \textbf{(e)} Justify with composition --- both orders --- that $f$ and $f^{-1}$ are inverses.
        \textbf{(f)} A classmate writes $y=(x-1)^2+4$ with \emph{no} restriction and calls it the
        inverse. Explain why that is not the inverse of $f$.
        \vspace{6.2cm}

  \item \textbf{Modeling.} Highway engineers estimate the greatest safe speed on a flat curved ramp
        of radius $r$ feet as
        \[ s(r)=4\sqrt{r}\ \text{ miles per hour.} \]
        \textbf{(a)} State the domain that makes sense here, and connect it to the rule about what an
        even index will accept.
        \textbf{(b)} Find $s(25)$ and $s(100)$, with units. What does quadrupling the radius do to the
        safe speed, and which feature of $\sqrt{\,\cdot\,}$ is responsible?
        \textbf{(c)} A ramp is posted at $36$ miles per hour. Find its radius algebraically, and check
        your answer.
        \textbf{(d)} Write the inverse, $r$ as a function of $s$, and say what it tells an engineer
        that $s(r)$ does not.
        \textbf{(e)} Part (c) required squaring both sides. Explain why there is no danger of an
        extraneous solution in this problem.
        \vspace{5.5cm}
\end{enumerate}

\end{document}
